\section{Engel Theorem the first type}

\begin{definition}
    Let \(L\) be a Lie algebra over a field \(F\). An \emph{enveloping algebra} of \(L\) is an associative \(F\)-algebra \(A\) together with a Lie algebra homomorphism
    \[
        \iota:L\to A^{-},
    \]
    where \(A^{-}\) denotes the Lie algebra obtained from \(A\) by equipping the underlying vector space of \(A\) with the bracket
    \[
        [x,y]=xy-yx,
    \]
    such that the associative algebra \(A\) is generated by \(\iota(L)\) (or just writing $A=<L>_{assoc}$ or simply $A=<L>$).

    If, in addition, the pair \((A,\iota)\) satisfies the following universal property: for every associative \(F\)-algebra \(B\) and every Lie algebra homomorphism
    \[
        \varphi:L\to B^{-},
    \]
    there exists a unique associative algebra homomorphism
    \[
        \widetilde{\varphi}:A\to B
    \]
    such that
    \[
        \varphi=\widetilde{\varphi}\circ \iota,
    \]
    then \(A\) is called the \emph{universal enveloping algebra} of \(L\).
\end{definition}

\begin{theorem}
    Let $ A $ be a finite dimensional enveloping algebra of $L$, i.e. $ L\subseteq A^{(-)} $, $ A = <L> $.If $ \exists n \ge 1 $ such that $ \forall a \in L $, $ a^n=0 $. Then $ A $ is nilpotent.
\end{theorem}

